\documentstyle[%


righttag%


,11pt%


,amssymb%


,russian%               remove in Eng.


,izvru%                 Set izven% in Eng.


]{amsart}





% ! {=} {\in}


\newcommand{\sign}{\operatorname{sign}}


\newcommand{\grad}{\operatorname{grad}}


%%%%%%%%%    vvvvvvvv               extremely widely used


\newcommand{\op}{{\text{оп}}} %% !! take account of in Eng.


\newcommand{\nop}{{\text{н}}} %% !!


\newcommand{\opn}{{\text{он}}} %% !!





\newcommand{\co}{\operatorname{co}}


\newcommand{\la}{\langle}


\newcommand{\ra}{\rangle}


\newcommand{\tts}[1]{}





\theoremstyle{plain}


\theorembodyfont{\it}


\newtheorem{thmnonum}{\hskip\parindent Теорема}


\renewcommand{\thethmnonum}{}





\theoremstyle{definition}


\theorembodyfont{\rm}


\newtheorem{notenonum}{\hskip\parindent Замечание}


\renewcommand{\thenotenonum}{}





%\hoffset=-15mm%                  For Eng.


\setcounter{page}{3}


\god{2004}


\nomer{12~(511)} %                        Remove in Eng.


%\nomer{Vol.\,48, No.\,12}%               Set in Eng.


%\str{\pageref{begin}--\pageref{end}}%  Set in Eng.


\udk{517.977}





\author{\it Р.\,Габасов, Ф.М.\,Кириллова, Т.Г.\,Хомицкая}


% NB:   ^^ remove in Eng.


\title{Программное и позиционное решения терминальной


линейно выпуклой задачи оптимального управления}


\thanks{Работа выполнена при финансовой поддержке Белорусского


Республиканского фонда


фундаментальных исследований (проект \No\,Ф04Р-002) и Государственной


программы фундаментальных исследований (``Математические


структуры--16'').}





\date{}


%\pagestyle{myheadings}%         Set in Eng.


\pagestyle{plain} %              Remove in Eng.


\begin{document}


%\thispagestyle{plain}%          Set in Eng.


\maketitle


\label{begin}











\section{Введение}





Линейно выпуклые задачи оптимального управления (ОУ) являются 


непосредственным обобщением линейных


задач. В приложениях они часто встречаются


как самостоятельные, в конструктивной теории могут использоваться и как


вспомогательные при решении нелинейных


задач ОУ [1]. Качественная теория таких задач разработана столь же глубоко


[2], как и аналогичная теория для


линейных задач. Предложено много алгоритмов их численного решения. Однако


известные алгоритмы касаются в основном


программной оптимизации динамических систем и не настолько эффективны,


чтобы синтезировать ОУ типа обратной связи.


Исключение составляют линейно-квадратичные задачи Летова--Калмана и


примыкающие к ним задачи $H_\infty$-теории


управления, для которых удается построить ОУ в виде {\it линейных}


обратных связей. Эти результаты привлекли


большое внимание в силу простоты их реализации. Задачи Летова--Калмана


являются гладкими и относятся скорее к


классическому вариационному исчислению, чем к теории ОУ, поскольку в них


игнорируются прямые (геометрические)


ограничения на управления. Подобные ограничения делают экстремальные


задачи негладкими и представляют наиболее


распространенный (нелинейный) элемент современных прикладных задач. Именно


из-за них в свое время пришлось строить


теорию вариационных задач неклассического типа, которая стала называться


теорией ОУ. За пятьдесят лет теория ОУ


достигла выдающихся успехов. Однако в ней до сих пор остается нерешенной


центральная проблема --- синтез ОУ типа


обратной связи


--- даже для линейных систем.





Частный случай исследуемых в работе задач --- линейно-квадратичные задачи


ОУ --- играют большую роль в бурно


развивающейся теории стабилизации, основанной на методологии Model


Predictive Control (MPC) [3]. Прикладное


значение этой теории особенно возросло [4], когда в ней стали


рассматриваться задачи ОУ с прямыми ограничениями на


управления. Реализуемость методологии MPC зависит от возможности


построения в режиме реального времени программных


решений специальных задач оптимального управления. Только в этом случае


можно в конкретных процессах стабилизации


формировать управляющие воздействия по ходу процесса стабилизации


(on-line control). В известной литературе по


MPC пока нет специальных быстрых методов вычисления соответствующих


программных решений задач ОУ. Для этой цели там


используют в основном стандартные методы квадратичного программирования,


что ограничивает возможности MPC


медленными переходными процессами из-за необходимости использования


больших периодов дискретизации времени и малых


горизонтов управления. В данной работе описываются специальные методы


решения линейно выпуклых задач ОУ,


эффективность которых мало зависит от периода дискретизации времени.


Первые методы подобного типа для линейных


задач ОУ были предложены еще в [5], [6], затем они получили развитие в


[7]--[9]. Результаты по ОУ в режиме


реального времени были независимо от MPC использованы в [10]--[13] для


стабилизации динамических систем


ограниченными управлениями.





{\it Цель данной работы} --- построить методы программной и позиционной


оптимизации линейных нестационарных


систем по выпуклым терминальным критериям качества.





Предлагается метод вычисления оптимальной программы


кусочно-линейной аппроксимации исходной задачи.


Для решения


линейных задач ОУ с параметром разработан быстрый двойственный


метод и процедура доводки.


Oбсуждается вопрос о вычислении оптимальных программ в классе


дискретных управлений.


Вычисление оптимальных программ не является конечной целью


конструктивной теории ОУ. Оптимальные программы выявляют лишь


потенциальные возможности систем управления, но они, как правило,


не участвуют в процессах реального управления. Реальные процессы


управления базируются на управлениях типа обратной связи. Проблема


синтеза оптимальных систем в классической постановке до сих пор


остается нерешенной. Данная работа основана на другом подходе к


проблеме [5]. При этом подходе описывается метод


реализации оптимальных обратных связей в режиме реального времени.








\section{Линейно выпуклая задача терминального управления.\protect\\


Общая схема решения задачи}





Пусть $T=[t_*,t^*]$, $t_*<t^*<+\infty$, --- промежуток управления,


$A(t)$, $b(t)$, $t\in T$, --- кусочно-непрерывные $n\times


n$-матричная и $n$-векторная функции, $\varphi(x)$, $x\in R^n$,


в отличие от [8] --- достаточно гладкая выпуклая функция:


$\varphi(x)\rightarrow\infty$ при $\|x-x^*\|\rightarrow\infty$,


$x^*$ --- точка минимума:


$\varphi^*=\varphi(x^*)=\min\,\varphi(x)$, $x\in R^n$.





В классе скалярных кусочно-непрерывных управляющих воздействий


$u(t)$, $t\in T$, рассмотрим линейно выпуклую задачу ОУ


\begin{equation}\label{problem}


\alpha^0=\min\varphi(x(t^*)),\ \ \Dot{x}=A(t)x+b(t)u,


\ \ x(t_*)=x_0,\ \ |u(t)|\leq1,\ \ t\in T.


\end{equation}


Здесь $x=x(t)$ --- $n$-вектор состояния системы управления в


момент времени $t$, $u=u(t)$ --- значение скалярного управляющего


воздействия, $x_0\in R^n$ --- заданное начальное состояние.





Кусочно-непрерывное управление $u^0(t)$, $|u^0(t)|\leq1$, $t\in


T$, назовем оптимальной программой, если на соответствующей ей


траектории $x^0(t)$, $t\in T$, критерий качества задачи


(\ref{problem}) достигает минимального значения:


$\alpha^0=\varphi(x^0(t^*))=\min\varphi(x(t^*))$, где минимум


берется по всем кусочно-непрерывным программам $|u(t)|\leq 1$,


$t\in T$.





Для решения поставленной задачи программной оптимизации линейной


нестационарной системы (\ref{problem}) в классе


кусочно-непрерывных управлений введем кусочно-линейную задачу


оптимизации в классе дискретных программ.





Управление $u(t)$, $t\in T$, назовем дискретной программой (с


периодом квантования $h$, $h=(t^*-t_*)/N$, $N$


--- натуральное число), если $u(t)=u(t_k)$,


$t\in [t_k,t_{k+1}[$, $t_k=t_*+kh$, $k=\overline{0,N-1}$.





Пусть $X^\alpha=\{x\in R^n:\varphi(x)\leq\alpha\}$, $\alpha\in


R$. Задачу (\ref{problem}) запишем в эквивалентном виде


\begin{equation}\label{ekviv-problem}


\alpha^0=\min\alpha,\ \ \Dot{x}=A(t)x+b(t)u,


\ \ x(t_*)=x_0,\ \ x(t^*)\in X^\alpha,\ \ |u(t)|\leq1,\ \ t\in T.


\end{equation}





Выберем произвольную совокупность единичных $n$-векторов


\begin{equation}\label{sovokupnost}


h_i,\ \ i\in I=\{1,2,\dotsc,m\},\ \ m\geq n,


\end{equation}


в которой каждые $n$ векторов линейно независимы. Векторы


(\ref{sovokupnost}) назовем векторами аппроксимации. В качестве


таких векторов можно взять единичные орты


$e_i$, $i=\overline{1,n}$.





Обозначим $r_i=h_i$, $r_{m+i}=-h_i$, $i=\overline{1,m}$, и


подсчитаем числа $g^*_i(\alpha)=\max r'_ix$, $x\in X^\alpha$,


${g_*}_i(\alpha)=\max r'_{m+i}x$, $x\in X^\alpha$, $i\in I$.





Множество $\wt{X}^\alpha=\{x\in R^n:{g_*}_i(\alpha)\leq


h'_ix\leq g^*_i(\alpha)$, $i\in I\}$ называется (внешней)


кусочно-линейной аппроксимацией множества $X^\alpha$.





Под кусочно-линейной аппроксимацией задачи (\ref{ekviv-problem})


понимается следующая задача ОУ в классе дискретных управлений:


\begin{equation}\label{approxim-problem}


\wt{\alpha}^0=\min\alpha,\ \ \Dot{x}=A(t)x+b(t)u,


\ \ x(t_*)=x_0,\ \ x(t^*)\in \wt{X}^\alpha,\ \ |u(t)|\leq1,\ \ t\in T.


\end{equation}





Оптимальные программы задачи (\ref{problem}) будем строить с


помощью двух процедур:





\noindent{}1)~вычисление оптимальной программы


кусочно-линейной задачи (\ref{approxim-problem});


2)~процедура доводки.





\noindent{}Первая процедура будет описана в п.\,3, вторая --- в


п.\,4.





Выбрав период квантования $h$ и совокупность (\ref{sovokupnost}),


вычислим программное решение $u_{\alpha}^0(t)$, $x_{\alpha}^0(t)$,


$t\in T$, задачи (\ref{approxim-problem}). Подсчитаем


$\varepsilon_i=g_i^*(\alpha)-r_i'x_\alpha^0(t^*)$,


$\varepsilon_{m+i}=-{g_*}_i(\alpha)-r_{m+i}'x_\alpha^0(t^*)$,


$i\in I$. Среди чисел $\varepsilon_i$, $i=\overline{1,2m}$,


выберем $n$ наименьших: $\varepsilon_i$, $ i\in I'$, $|I'|=n$.


Число $\wt{\varepsilon}=\max\varepsilon_i$, $i\in I'$,


характеризует точность аппроксимaции программного решения исходной


задачи программным решением задачи (\ref{approxim-problem}). Если


при


выбранном\footnote{Число $\varepsilon$ подбирается так,


чтобы последующая процедура доводки сходилась.}


$\varepsilon>0$


выполняется неравенство $\wt{\varepsilon}\leq\varepsilon$, то


перейдем к процедуре доводки (п.\,4). При


$\wt{\varepsilon}>\varepsilon$ к векторам $r_i$, $i\in I$,


добавим вектор


$\wt{r}=\grad\varphi(x_\alpha^0(t^*))/\|\grad


\varphi(x_\alpha^0(t^*))\|$. Вычислим


$\wt{g}^*(\alpha)=\max\wt{r}'x$, $x\in X^\alpha$;


$\wt{g}_*(\alpha)=\min\wt{r}'x$, $x\in X^\alpha$, и в


задаче (\ref{approxim-problem}) заменим терминальные ограничения


на следующие:


$$


{g_*}_i(\alpha)\leq r'_ix(t^*)\leq g^*_i(\alpha),\ \ i\in I;


\ \ \wt{g}_*(\alpha)\leq \wt{r}'x(t^*)\leq \wt{g}^*(\alpha).


$$


Программное решение новой задачи можно быстро построить с помощью


двойственного метода, описанного в п.\,3. При дальнейшем повышении


точности аппроксимации задача (\ref{approxim-problem}) решается с


новыми терминальными ограничениями, которые получаются из


предыдущих добавлением одного нового ограничения, построенного по


приведенным выше правилам. Когда количество новых ограничений


достигнет $n$, удалим из добавленных одно (``худшее'', с


наибольшим $\varepsilon_i$) ограничение. Через конечное число


коррекций терминальных ограничений получается программное решение


задачи (\ref{approxim-problem}), на котором будут выполняться


условия перехода к процедуре доводки. Процедура доводки строит


программное решение задачи (\ref{problem}) в классе


кусочно-непрерывных функций.








\section{Вычисление оптимальной программы кусочно-линейной задачи ОУ}





Задача (\ref{approxim-problem}) является нелинейной задачей ОУ.


Решим ее с помощью конечного числа коррекций терминальных


ограничений линейной задачи. Положив $l=x^*-x_0$, для


фиксированного $\alpha$ в классе дискретных программ рассмотрим


линейную задачу ОУ


\begin{gather}\label{vspom-problem}


\beta(\alpha)=\min\beta,\ \ \Dot{x}=A(t)x+b(t)u,


\ \ x(t_*)=x_\beta=x_0+\beta l,\\


x(t^*)\in \wt{X}^\alpha,\ \ |u(t)|\leq1,\ \ t\in T\ \ (\beta\in R).\notag


\end{gather}





Для $\alpha=\varphi^*$ задача (\ref{vspom-problem}) имеет вид


\begin{gather}\label{vspom-problem-inval}


\beta^*=\min\beta,\ \ \Dot{x}=A(t)x+b(t)u,


\ \ x(t_*)=x_\beta=x_0+\beta l,\\


x(t^*)=x^*,\ \ |u(t)|\leq1,\ \ t\in T.\notag


\end{gather}





Нетрудно показать справедливость следующих утверждений.





\begin{lem}


Функция $\beta(\alpha)$, $\alpha\in R$, является непрерывной


монотонно убывающей функцией\/$;$


$\beta(\varphi^*)=\beta^*$, $\beta(\alpha)\rightarrow-\infty$ при


$\alpha\rightarrow\infty$.


\end{lem}





\begin{lem}


Оптимальное значение $\wt{\alpha}^0$ критерия


качества задачи $(\ref{approxim-problem})$


равно минимальному корню уравнения $\beta(\alpha)=0$.


\end{lem}





Двойственный метод решения задачи (\ref{vspom-problem}) приведен


ниже.





Решение задачи (\ref{approxim-problem}) начнем с решения задачи


(\ref{vspom-problem-inval}). Если $\beta^*=\beta(\varphi^*)<0$, то


$x^*$ --- внутренняя точка множества достижимости системы


(\ref{problem}), в задаче (\ref{approxim-problem}) заведомо


существует бесконечное множество решений, на которых можно ввести


дополнительный критерий качества. В случае $\beta^*=0$ (согласно


лемме 2) $\varphi^*$ --- минимальное значение критерия качества


задачи (\ref{approxim-problem}): $\wt{\alpha}^0=\varphi^*$. При


$\beta^*>0$ найдем минимальный корень уравнения $\beta(\alpha)=0$.


Это уравнение можно решать любым из многочисленных известных


методов.





Изложим двойственный метод решения задачи


(\ref{vspom-problem}). Сначала приведем необходимые элементы этого


метода [14] в форме, удобной для решения исследуемой задачи.





Пусть $T_h=\{t_*,t_*+h,\dotsc,t^*-h\}$. Используя формулу Коши [2]


$x(t^*)=F(t^*,t_*)x(t_*)+


\int\limits_{t_*}^{t^*}F(t^*,\tau)b(\tau)u(\tau)d\tau=


F(t^*,t_*)x_0+F(t^*,t_*)\beta l+\sum\limits_{t\in T_h}


\int\limits_t^{t+h}F(t^*,\tau)b(\tau)d\tau\,u(t)$, где


$F(t,\tau)=F(t)F^{-1}(\tau)$, $\Dot{F}=A(t)F$, $F(t_*)=E$, запишем


задачу (\ref{vspom-problem}) в эквивалентной функциональной форме


\begin{equation}\label{funkform-problem}


\beta(\alpha)=\min\beta,\ \ q_*(\alpha)\leq\sum\limits_{t\in


T_h}f_h(t)\,u(t)+f_l\beta\leq q^*(\alpha),\ \ |u(t)|\leq1,\ \ t\in T.


\end{equation}


Здесь


\begin{gather}\label{element-problem}


f_h(t)=\int_t^{t+h}HF(t^*,\tau)b(\tau)d\tau,\ \ t\in


T_h,\ \ f_l=HF(t^*,t_*)l,\\


q_*(\alpha)=g_*(\alpha)-HF(t^*,t_*)x_0,


\ \ q^*(\alpha)=g^*(\alpha)-HF(t^*,t_*)x_0,\notag


\end{gather}


$H$ --- $m\times n$-матрица, составленная из строк $h_i$, $i\in I$.





Динамический способ построения элементов (\ref{element-problem})


состоит в следующем. Пусть $G(t)$, $t\in T$, --- $m\times


n$-матричная функция --- решение уравнения


\begin{eqnarray}\label{dif-uravn-G}


 & \dot{G}=-GA(t),\ \ G(t^*)=H.


\end{eqnarray}


Тогда элементы (\ref{element-problem}) можно строить по формулам


$$


f_h(t)=\int_t^{t+h}G(\tau)b(\tau)d\tau,


\ \ f_l=G(t_*)l,\ \ q_*(\alpha)=g_*(\alpha)-G(t_*)x_0,


\ \ q^*(\alpha)=g^*(\alpha)-G(t_*)x_0.


$$





Задача (\ref{funkform-problem}) отличается от соответствующей


задачи из [15] наличием параметра $\beta$, что требует внесения


некоторых дополнений в конструкции [15].





Начнем с основного понятия метода --- опоры. Выделим из множества


индексов $I$ произвольное подмножество $I_{\op}$, $|I_{\op}|=p$,


$1\leq p\leq m$, а из $T_h$


--- произвольное подмножество $T_{\op}=\{t_j$, $j=\overline{1,p-1}\}$


и построим $p\times p$-матрицу


$P_{\op}=\left(\begin{smallmatrix}


f_{hi}(t),\ t\in T_{\op}; \ f_{li} \\


i\in I_{\op}


\end{smallmatrix}\right)$, где $f_{hi}(t)$, $f_{li}$


--- $i$-е координаты векторов $f_h(t)$, $f_l$.





\begin{defn}


Пару множеств $K_{\op}{=}\{I_{\op},T_{\op}\}$


назовем опорой задачи (\ref{funkform-problem}),


если $\det P_{\op}\neq0$.


\end{defn}





Опору, составленную из $I_{\op}=\{i_1: f_{li_1}\neq0\}$ и


$T_{\op}=\emptyset$, будем называть простейшей опорой.





Опишем два способа динамической проверки совокупности $K_{\op}$ на


опорность.





{\it Прямой способ.}





1) Построим терминальные состояния $\chi_\tau(t^*)$, $\tau\in


T_{\op}$, прямой системы (\ref{problem}), соответствующие


начальному состоянию $x(t_*)=0$ и программе


$$


u(t)=1,\ \ t\in [\tau,\tau+h[;\ \ u(t)=0,


\ \ t\in T\setminus [\tau,\tau+h[.


$$





2) Найдем терминальное состояние $\chi^\beta(t^*)$ прямой системы


(\ref{problem}), соответствующее начальному состоянию $x(t_*)=l$ и


программе $u(t)=0$, $t\in T$.





3) Умножив матрицу $H_{\op}=(h_i,\ i\in I_{\op})'$ на эти


состояния, составим матрицу


\begin{equation}\label{matrix-P-1}


P_{\op}=(H_{\op}\chi_{\tau}(t^*),\ \tau\in


T_{\op};\ H_{\op}\chi^\beta(t^*)).


\end{equation}





{\it Двойственный способ.}





1) Проинтегрировав сопряженную систему


\begin{equation}\label{sopr-sistem}


\dot{\psi}=-A'(t)\psi


\end{equation}


с начальными условиями $\psi(t^*)=h_i$, $i\in I_{\op}$, найдем


\begin{equation*}


\xi'_i(t)=h'_iF(t^*,t),\ \ t\in T,\ \ i\in I_{\op}.


\end{equation*}





2) Составим матрицу


\begin{equation}\label{matrix-P-2}


P_{\op}=\begin{pmatrix}


\int\limits_t^{t+h}\xi'_i(\tau)b(\tau)d\tau,\ t\in T_{\op}; \ \ 


\xi'_i(t_*)l \\


i\in I_{\op}


\end{pmatrix}.


\end{equation}





Таким образом, совокупность $K_{\op}$ является опорой тогда и


только тогда, когда любая из матриц (\ref{matrix-P-1}),


(\ref{matrix-P-2}) неособая.





Опору $K_{\op}$ сопровождают следующие элементы.





1. Вектор потенциалов $\nu=(\nu_1,\dotsc, \nu_p)$. Вычислим его


как решение уравнения $\nu'P_{\op}=e'_p$, где $e_p=(0, 0,\dotsc,


1)\in R^p$


--- единичный орт. Ясно, что $\nu'=\big[P_{\op}^{-1}\big]_p$, т.\,е. $\nu$


представляет $p$-ю строку матрицы $P_{\op}^{-1}$.





2. Котраектория $\psi(t)$, $t\in T$, --- решение сопряженного


уравнения (\ref{sopr-sistem}) с начальным условием


$\psi(t^*)=-H'_{\op}\nu$.





3. Коуправление $\delta_h(t)$, $t\in T_h$, --- динамический аналог


вектора оценок [16]: $\delta_h(t)=-\sum\limits_{i\in


I_{\op}}\nu_i{f_h}_i(t)$, $t\in T_h$. Динамическая формула для


вычисления коуправления имеет вид


\begin{equation*}


\delta_h(t)=\int_t^{t+h}\psi'(\tau)b(\tau)d\tau,\quad


t\in T_h.


\end{equation*}





4. Псевдоуправление $\omega(t)$, $t\in T_h$, параметр


$\beta_\omega$ начального состояния и выходной псевдосигнал


$\zeta(t^*)=(\zeta_i(t^*),\ i\in I)$ в момент $t^*$. Для их


построения сначала зададим опорные элементы


\begin{gather*}


\zeta_i(t^*)={g_*}_i(\alpha) \text{ при }\nu_i>0;\quad


\zeta_i(t^*)=g^*_i(\alpha) \text{ при }\nu_i<0;\\


\zeta_i(t^*)\in [{g_*}_i(\alpha);g^*_i(\alpha)]


\text{ при }\nu_i=0,\ \ i\in I_{\op},


\end{gather*}


и значения псевдоуправления в неопорные моменты


\begin{gather*}


\omega(t)=-1 \text{ при }\delta_h(t)>0;\quad


\omega(t)=1 \text{ при }\delta_h(t)<0;\\


\omega(t)\in [-1; 1] \text{ при }\delta_h(t)=0,\,t\in T_{\nop}.


\end{gather*}





Динамический способ построения совокупности


$\{\omega_{\op}=(\omega(t),\ t\in T_{\op})$; $\beta_{\omega}\}$


сводится к следующим операциям.





1)~Вычислим $\varkappa_0(t^*)$ --- значение в момент $t^*$ решения


$\varkappa_0(t)$, $t\in T$, прямого уравнения (\ref{problem}) с


начальным условием $x(t_*)=x_0$ и программой $u(t)=\omega(t)$,


$t\in T_{\nop}$; $u(t)=0$, $t\in T_{\op}$.





2)~Решим уравнение


\begin{equation}\label{uravn-psevdoupr}


P_{\op}\begin{pmatrix}


\omega_{\op} \\ \beta_\omega


\end{pmatrix}=\zeta_{\op}-H_{\op}\varkappa_0(t^*),


\end{equation}


где $\zeta_{\op}=(\zeta_i(t^*),\ i\in I_{\op})$.





Решение $\varkappa(t)$, $t\in T$, прямого уравнения (\ref{problem}) с


начальным условием $x(t_*)=x_0+\beta_{\omega}l$ и дискретной


программой $u(t)=\omega(t)$, $t\in T$, назовем псевдотраекторией,


вектор $\zeta=\zeta(t^*)=H\varkappa(t^*)$ --- выходным псевдосигналом в


момент $t^*$, сопровождающими опору $K_{\op}$.





Неопорные компоненты вектора $\zeta(t^*)$ равны


$\zeta_{\nop}=(\zeta_i(t^*),\ i\in I_{\nop})=(h'_i\varkappa(t^*),\ i\in


I_{\nop})$.





Опору используем прежде всего для идентификации оптимальных


программ.





\begin{thm}[{\series{m}\shape{n}\selectfont{}принцип минимума}]


Для оптимальности программы


$u(t)$, $t\in T$, необходимо и достаточно


существования такой опоры $K_{\text{\em оп}}$, %%% <-- !


чтобы на сопровождающих ее векторе


потенциалов $\nu$ и котраектории $\psi(t)$,


$t\in T$, выполнялись два условия\/${:}$





$1)$ условие минимума для программы


\begin{equation}\label{krit-optimaln:1}


\int_t^{t+h}\psi'(\tau)b(\tau)d\tau\,u(t)=\min\limits_{|u|\leq1}


\int_t^{t+h}\psi'(\tau)b(\tau)d\tau\,u,\quad  t\in T_h,


\end{equation}





$2)$ условие трансверсальности для траектории


\begin{equation}\label{krit-optimaln:2}


\nu'Hx(t^*)=\min \nu'z,\ \ g_*(\alpha)\leq z\leq g^*(\alpha).


\end{equation}


\end{thm}





\begin{defn}


Опору $K_{\op}$, на которой выполняются условия


(\ref{krit-optimaln:1}),


(\ref{krit-optimaln:2}), назовем оптимальной.


\end{defn}





\begin{thm}[{\series{m}\shape{n}\selectfont{}критерий оптимальности опоры}]


Для оптимальности опоры $K_{\text{\em оп}}$ %%% <-- !


необходимо и достаточно,


чтобы на некоторых сопровожающих ее псевдоуправлении $\omega(t)$, $t\in


T$, и выходном псевдосигнале


$\zeta(t^*)=H\varkappa(t^*)$ выполнялись следующие неравенства\/${:}$


\begin{equation*}


|\omega(t)|\leq1,\ \ t\in T_{\text{\em оп}};


\ \ {g_*}_i(\alpha)\leq\zeta_i(t^*)


\leq g^*_i(\alpha),\ \ i\in I_{\text{\em н}}. %%% <-- !


\end{equation*}


\end{thm}





\begin{defn}


Опору $K_{\op}$ назовем регулярной, если


сопровождающие ее вектор потенциалов $\nu$ и


коуправление $\delta_h(t)$, $t\in T_h$, удовлетворяют соотношениям


$\nu_i\neq0$, $i\in I_{\op}$;


$\delta_h(t)\neq0$, $t\in T_{\nop}$.


\end{defn}





Регулярную опору сопровождают единственные псевдоуправление


$\omega(t)$, $t\in T$, и выходной псевдосигнал $\zeta(t^*)$.





На основе этих результатов разработан {\it двойственный метод}


решения линейной задачи (\ref{vspom-problem}), который


представляет последовательное преобразование опор до получения


оптимальной опоры.





Пусть $K_{\op}$ --- начальная опора задачи (в качестве ее можно взять


простейшую опору). Предположим, что она не


удовлетворяет критерию оптимальности опоры.





С целью упрощения выкладок предположим, что 1)~на итерациях метода


участвуют лишь регулярные опоры (общий случай


рассмотрен в [16]); 2)~коуправление $\delta_h(t)$, $t\in T_h$,


удовлетворяет условиям


\begin{gather*}


\delta_h(t-h)\delta_h(t+h)<0, \text{ если } t\in T_{\op},


\ \ t_*<t<t^*-h;\\


\delta_h(t_*+h)\neq0, \text{ если } t^*\in T_{\op};\ \


\delta_h(t_*-2h)\neq0, \text{ если } t^*-h\in T_{\op}.


\end{gather*}





Момент $t\in T_{\nop}\setminus t_*$ будем называть неопорным нулем, если


$\delta_h(t-h)\delta_h(t)<$~$0$.


Множество неопорных нулей обозначим через $T_{\nop 0}$. Пусть


$T_{\opn}=T_{\nop 0}\cup T_{\op}\cup


\{t_*,t^*\}=\{t_s,\ s\in S\cup s^*+1\}$, $S=\{0, 1,\dotsc, s^*\} $.


Через $T_s$, $s\in S$, обозначим промежутки


постоянства знака коуправления:


\begin{alignat*}{2}


T_s&=\{{t_*}_s=t_s, t_s+h,\dotsc, t_s^*=t_{s+1}-h\},&\quad &\text{если }


t_s\notin T_{\op};\\


T_s&=\{{t_*}_s=t_s+h, t_s+2h,\dotsc, t_s^*=t_{s+1}-h\},&\quad &\text{если }


t_s\in T_{\op}.


\end{alignat*}


Если $t^*-h\in T_{\op}$, то положим $T_{s^*}=0$.





Введем векторы


\begin{equation*}


d=H\varkappa_0(t^*)=G(t_*)x_0+\gamma\sum\limits_{s=0}^{s^*}(-


1)^{s+1}\int_{{t_*}_s}^{t_s^*+h}G(\tau)b(\tau)d\tau,\quad


d_{\op}=(d_i,\ i\in I_{\op})


\end{equation*}


($\gamma=\sign\delta_h(t_*)$, если $t_*\notin T_{\op}$;


$\gamma=\sign\delta_h(t_*+h)$, если $t_*\in T_{\op}$).





Тогда уравнение (\ref{uravn-psevdoupr}) примет вид


$P_{\op}\left(\begin{smallmatrix}


\omega_{\op} \\ \beta_\omega


\end{smallmatrix}\right)=\zeta_{\op}-d_{\op}$.





С помощью вектора $d$ и матрицы $P_{|\op|}=(f_h(t),\ t\in T_{\op};\ f_l)$


легко подсчитать и выходной псевдосигнал


$\zeta(t^*)=P_{|\op|}\left(\begin{smallmatrix}


\omega_{\op} \\ \beta_\omega


\end{smallmatrix}\right)+d$. Матрица $P_{|\op|}$ строится аналогично


$P_{\op}$, например, прямым способом:


$P_{|\op|}=(H\chi_{\tau}(t^*)$, $\tau\in T_{\op}$; $H\chi^{\beta}(t^*))$ или


двойственным: $P_{|\op|}=\Big(\int\limits_t^{t+h}G(\tau)b(\tau)d\tau$,


$t\in T_{\op}$; $G(t_*)l\Big)$.





Будем считать, что к началу итерации кроме параметров задачи $A(t)$,


$b(t)$, $t\in T$; $H$, $g_*(\alpha)$,


$g^*(\alpha)$, $l$, известна следующая информация (хранится в оперативной


памяти ЭВМ): 1)~опора


$K_{\op}=\{I_{\op},T_{\op}\}$, 2)~множество неопорных нулей $T_{\nop0}$,


3)~матрица $P_{|\op|}$, 4)~значения


$G(t)$, $t\in T_{\opn}\setminus t^*$, 5)~величина $\gamma$, 6)~опорные


значения псевдоуправления $\omega(t)$,


$t\in T_{\op}$, и значение параметра $\beta_{\omega}$, 7)~вектор $d$,


8)~вектор потенциалов $\nu$, 9)~вектор


выходного псевдосигнала $\zeta(t^*)$.





Информация, необходимая для первой итерации, готовится по изложенным выше


правилам. Далее она преобразуется по ходу


итераций.





{\it Итерация.} По компонентам $\omega(t)$, $t\in T_{\op}$;


$\zeta_i(t^*)$, $i\in I_{\nop}$, вычислим


\begin{equation}\label{rasstojanije}


\rho(\zeta_i(t^*),[{g_*}_i(\alpha),g^*_i(\alpha)]),\ \ i\in


I_{\nop};\ \; \rho(\omega(t),[-1,1]),\ \ t\in T_{\op},


\end{equation}


($\rho(c,\,[a,b])$ --- расстояние от точки $c$ до  отрезка $[a,b]$). Из


(\ref{rasstojanije}) выберем наибольшие


\begin{align*}


\rho_{i_0}=\rho(\zeta_{i_0}(t^*),[{g_*}_{i_0}(\alpha),g^*_{i_0}(\alpha)]


)&=\max\rho(\zeta_i(t^*),[{g_*}_i(\alpha),g^*_i(\alpha)]),\ \ i\in


I_{\nop};\\


\rho(t_0)=\rho(\omega(t_0),[-1,1])&=\max\rho(\omega(t),[-


1,1]),\ \ t\in T_{\op}.


\end{align*}


Определим $\rho_0=\max\{\rho_{i_0},\rho(t_0)\}$ и исследуем каждую из


двух ситуаций: 1)~$\rho_0=\rho(t_0)$, \label{rho}


2)~$\rho_0=\rho_{i_0}$.





1)~($\rho_0=\rho(t_0)$). Начнем с формирования направления $\Delta\nu$


изменения вектора потенциалов $\nu$, решив


систему линейных уравнений


\begin{equation*}


-P'_{\op}\Delta\nu=(\Delta\delta_h(t),\ t\in


T_{\op})\quad (\Delta\delta_h(t_0)=-


\sign\omega(t_0);\ \Delta\delta_h(t)=0,\ t\in T_{\op}\setminus t_0).


\end{equation*}





Обозначим $\Delta\widetilde{\nu}=(\Delta\widetilde{\nu}_i,\ i\in I)$,


$\Delta\widetilde{\nu}_i=\Delta\nu_i$, если


$i\in I_{\op}$; $\Delta\widetilde{\nu}_i=0$, если $i\not\in I_{\op}$.


Направление изменения коуправления


$\Delta\delta_h(t)$, $t\in T_h$, построим по формуле


\begin{equation}\label{var-koupravlen}


\Delta\delta_h(t)=-


\int_t^{t+h}\Delta\widetilde{\nu}'G(\tau)b(\tau)d\tau,\quad t\in


T_h.


\end{equation}





Введем возмущенные коуправление


\begin{equation}\label{vozm-koupravlen}


\delta_h(t,\sigma)=\delta_h(t)+\sigma\Delta\delta_h(t),\ \ t\in


T_h;\ \ \sigma\geq0,


\end{equation}


и вектор потенциалов


\begin{equation}\label{vozm-potencial}


\nu(\sigma)=\nu+\sigma\Delta\nu.


\end{equation}





Далее предполагается, что $\Delta\delta_h({t_*}_s)\Delta\delta_h(t_{s-


1}^*)>0$, если ${t_*}_s\in T_{\opn}\setminus


\{t_*,\,t^*\}$, $s=\overline{1,s^*}$.





Согласно [16] в малой правосторонней окрестности точки $\sigma=0$


двойственный критерий качества изменяется вдоль


направления $\Delta\nu$  со скоростью


$\eta^{(1)}=\rho_0>0$.





В терминах вектора потенциалов замена опоры связана с движением вдоль


$\Delta\nu$ до полной релаксации в точке


$\overline{\nu}=\nu+\sigma^*\Delta\nu$ кусочно-линейного по $\sigma\geq0$


двойственного критерия качества.


``Длинный'' двойственный шаг $\sigma^*$ вычислим с помощью нескольких


``коротких'' шагов.





Для вычисления ``коротких'' шагов, в которых изменение двойственного


критерия качества терпит излом, наряду с


упомянутой информацией, хранящейся в памяти ЭВМ, будем использовать


дополнительную. Последняя для построения


первого ``короткого'' шага $\sigma^{(1)}$ имеет вид  $\eta^{(1)}$,


\begin{equation}\label{information}


\sigma(t_0),\ \; \vartheta(t_0),\ \; s_0;\ \; \sigma_i,\ \; i\in I_{\op};


\ \; \sigma(t_*),\ \; \sigma(t^*);\ \; \sigma(t),


\ \; \vartheta(t),\ \; t\in T_{\nop0}.


\end{equation}


Поясним физический смысл чисел (\ref{information}). Каждое из чисел


$\sigma(t)$, $\sigma_i$ представляет значение


шага вдоль $\Delta\nu$, при котором или компонента


(\ref{vozm-koupravlen}), или компонента (\ref{vozm-potencial})


обращается в нуль, что вызывает скачок скорости двойственного критерия


качества. Числа $\vartheta(t)$ указывают


направление перемещения нуля коуправления при увеличении $\sigma$; $s_0$


--- индекс момента $t_0$ в $T_{\opn}$.





Отсюда следуют формулы для вычисления чисел (\ref{information})


\begin{gather*}


\sigma(t_0)=-


\delta_h(t_0+h)/\Delta\delta_h(t_0+h),\ \ \vartheta(t_0)=1,


\ \ \text{если}\ \ (-1)^{s_0+1}\gamma\Delta\delta_h(t_0)<0;\\


\sigma(t_0)=-\delta_h(t_0-h)/\Delta\delta_h(t_0-h),\ \ \vartheta(t_0)=-1,


\ \ \text{если}\ \ (-1)^{s_0+1}\gamma\Delta\delta_h(t_0)>0,\\


\sigma_i=


\begin{cases}


-\nu_i/\Delta\nu_i, & \text{если }\nu_i\Delta\nu_i<0;\\


\infty, & \text{если }\nu_i\Delta\nu_i\geq0,\ \;i\in I_{\op},


\end{cases}


\\


\sigma(t_*)=\begin{cases}


-\delta_h(t_*)/\Delta\delta_h(t_*), & \text{если }


\delta_h(t_*)\Delta\delta_h(t_*)<0;\\


\infty, & \text{если }\delta_h(t_*)\Delta\delta_h(t_*)\geq0,


\end{cases}


\\


\sigma(t^*)=\begin{cases}


-\delta_h(t^*-h)/\Delta\delta_h(t^*-h), & \text{если }\delta_h(t^*-


h)\Delta\delta_h(t^*-h)<0;\\


\infty, & \text{если }\delta_h(t^*-h)\Delta\delta_h(t^*-h)\geq0,


\end{cases}


\\


\sigma(t)=-\delta_h(t)/\Delta\delta_h(t),\ \;\vartheta(t)=1,


\text{ если }\delta_h(t)\Delta\delta_h(t)<0;\\


\sigma(t)=-\delta_h(t-h)/\Delta\delta_h(t-h),\ \;\vartheta(t)=-1,


\text{ если }\delta_h(t-h)\Delta\delta_h(t-h)<0,\ \;t\in T_{\nop0};\\


\delta_h(t)=-\widetilde{\nu}'f_h(t),\ \; \Delta\delta_h(t)=-


\Delta\widetilde{\nu}'f_h(t),\ \;f_h(t)=\int_t^{t+h}G(\tau)


b(\tau)d\tau,\\


\delta_h(t+h)=-\widetilde{\nu}'f_h(t+h), \ \; \Delta\delta_h(t+h)=-


\Delta\widetilde{\nu}'f_h(t+h),\ \;f_h(t+h)=\int_{t+h}^{t+2h}G(\tau)


b(\tau)d\tau,\\


\delta_h(t-h)=-\widetilde{\nu}'f_h(t-h),\ \;


\Delta\delta_h(t-h)=-\Delta\widetilde{\nu}'f_h(t-h),\ \;f_h(t-


h)=\int_{t-h}^{t}G(\tau)b(\tau)d\tau,


\end{gather*}


$\widetilde{\nu}=(\widetilde{\nu}_i,\ i\in I)$, $\widetilde{\nu}_i=\nu_i$,


если $i\in I_{\op}$;


$\widetilde{\nu}_i=0$, если $i\notin I_{\op}$.





Таким образом, для вычисления чисел (\ref{information}) достаточно взять в


качестве начальных состояний $G(t)$


значения, хранящиеся в памяти ЭВМ, и проинтегрировать уравнение


(\ref{dif-uravn-G}) на промежутке $[t,t+2h]$ или $[t-h,t]$.





Прежде чем перейти к выполнению ``коротких'' шагов, проведем ``нулевой''


шаг, на котором преобразуем хранящуюся в


памяти ЭВМ информацию, считая, что сделан малый шаг $\sigma>0$ вдоль


направления $\Delta\nu$.





\begin{itemize}


\item[1.] При $t_*<t_0<t^*-h$, $\vartheta(t_0)=1$ положим


$d^{(1)}=d+(-1)^{s_0}\gamma f_h(t_0)$; $T_{\nop0}^{(1)}=T_{\nop0}\cup


t_0+h$; вместо значения $G(t_0)$ в память занесем $G(t_0+h)$.


\item[2.] Если $t_*<t_0\leq t^*-h$, $\vartheta(t_0)=-1$, то


$d^{(1)}=d+(-1)^{s_0+1}\gamma f_h(t_0)$, $T_{\nop0}^{(1)}=T_{\nop0}\cup


t_0$.


\item[3.] В случае $t_0=t_*$, $\vartheta(t_0)=1$ положим $d^{(1)}=d+\gamma


f_h(t_0)$, $T_{\nop0}^{(1)}=T_{\nop0}\cup


t_0+h$, $\gamma^{(1)}=-\gamma$; ${s^*}^{(1)}=s^*+1$;


перенумеруем точки множества $T_{\opn}$.


\item[4.] При $t_0=t_*$, $\vartheta(t_0)=-1$ считаем $d^{(1)}=d-\gamma


f_h(t_0)$, $T_{\nop0}^{(1)}=T_{\nop0}$. Случай интерпретируется как


исчезновение опорного


нуля через левую границу множества $T_h$.


\item[5.] Если $t_0=t^*-h$, $\vartheta(t_0)=1$, то положим $d^{(1)}=d+(-


1)^{s_0}\gamma f_h(t_0)$,


$T_{\nop0}^{(1)}=T_{\nop0}$; ${s^*}^{(1)}=s^*-1$, из памяти удалим


значение $G(t_0)$. Случай трактуется


как исчезновение опорного нуля через правую границу множества $T_h$.


\end{itemize}


Остальная информация для первого шага переписывается без изменений.





Предположим, что вдоль направления $\Delta\nu$ вычислен $\ell-1$


``короткий'' шаг $\sigma^{(1)}, \sigma^{(2)},


\dotsc, \sigma^{(\ell-1)}$ и информация перед $\ell$-м ``коротким'' шагом


имеет вид:





1)~множества


$T_{\nop0}^{(\ell)}$, $T^{(\ell)}=T_{\nop0}^{(\ell)}\cup t_*\cup t^*$,





2)~числа $\sigma^{(\ell)}(t)$, $t\in


T^{(\ell)}$, $\sigma_i^{(\ell)}$, $i\in I_{\op}$,





3)~число $\gamma^{(\ell)}$,





4)~вектор $d^{(\ell)}$,





5)~значения $G(t)$, $t\in T_{\nop0}^{(\ell)}\cup t_0$,





6)~число ${s^*}^{(\ell)}$,





7)~скорость изменения двойственного критерия


качества $\eta^{(\ell)}$.





С целью упрощения выкладок предположим, что все числа


$\sigma^{(\ell)}(t)$, $t\in T^{(\ell)}$; $\sigma_i^{(\ell)}$,


$i\in I_{\op}$, различны за исключением, возможно, пары чисел


$\sigma^{(\ell)}(t)$, $\sigma^{(\ell)}(t+h)$ при


некоторых $t\in T_h\setminus(t^*-h)$. Случай совпадения других чисел в


(\ref{information}) рассмотрен в [16].





Исходя из полученной информации, подсчитаем $\ell$-й ``короткий'' шаг


\begin{equation}\label{el-step}


\sigma^{(\ell)}=\min\{\sigma(t^{(\ell)}),\sigma_{i^{(\ell)}}\},


\end{equation}


где


$\sigma(t^{(\ell)})=\min\sigma(t)$, $t\in T^{(\ell)}$;


$\sigma_{i^{(\ell)}}=\min\sigma_i$, $i\in I_{\op}$.





Это позволит указать промежуток $[\sigma^{(\ell-1)},\sigma^{(\ell)}]$, на


котором двойственный критерий качества


возрастает со скоростью $\eta^{(\ell)}$. В точке $\sigma=\sigma^{(\ell)}$


скорость изменения двойственного критерия


качества вдоль $\Delta\nu$ убывает на


\begin{equation*}


\Delta\eta^{(\ell)}=\begin{cases}


2|\Delta\delta_h(t^{(\ell)})|, & \text{если}


\ \; \sigma^{(\ell)}=\sigma(t^{(\ell)});\\


(g_{i^{(\ell)}}^*(\alpha)-


{g_*}_{i^{(\ell)}}(\alpha))|\Delta\nu_{i^{(\ell)}}|, &


\text{если}\ \; \sigma^{(\ell)}=\sigma_{i^{(\ell)}},


\end{cases}


\end{equation*}


и в малой правосторонней окрестности точки $\sigma^{(\ell)}$ становится


равной


\begin{equation*}


\eta^{(\ell+1)}=\eta^{(\ell)}-\Delta\eta^{(\ell)}<\eta^{(\ell)}.


\end{equation*}


При $\eta^{(\ell+1)}\leq0$ положим $\sigma^*=\sigma^{(\ell)}$ и перейдем к


заключительному шагу. В противном случае


приступим к формированию информации 1)--7) для осуществления следующего


``короткого'' шага.





Будем различать две возможности: а)~$\sigma^{(\ell)}=\sigma(t^{(\ell)})$,


б)~$\sigma^{(\ell)}=\sigma_{i^{(\ell)}}$.





а) Хранящуюся в памяти информацию преобразуем в зависимости от следующих


ситуаций.





\noindent{}A. Пусть $\sigma(t^{(\ell)})$ --- единственное число,


удовлетворяющее (\ref{el-step}),


$s^{(\ell)}$ --- индекс момента $t^{(\ell)}$ в $T_{\opn}$.


\begin{itemize}


\item[A.1.] Если $t^{(\ell)}\in T_{\nop0}^{(\ell)}$ и


$\vartheta(t^{(\ell)})=1$,


 то положим $T_{\nop0}^{(\ell+1)}=(T_{\nop0}^{(\ell)}\setminus


t^{(\ell)})\cup t^{(\ell)}+h$;


 $d^{(\ell+1)}=d^{(\ell)}+2(-


1)^{s^{(\ell)}}\gamma^{(\ell)}f_h(t^{(\ell)})$; вычислим новый шаг


 $\sigma(t^{(\ell)}+h)=-


\delta_h(t^{(\ell)}+h)/\Delta\delta_h(t^{(\ell)}+h)$. Вместо значения


 $G(t^{(\ell)})$ в память ЭВМ занесем $G(t^{(\ell)}+h)$.


\item[A.2.] Если $t^{(\ell)}\in T_{\nop0}^{(\ell)}$ и


$\vartheta(t^{(\ell)})=-1$, то


 положим $T_{\nop0}^{(\ell+1)}=(T_{\nop0}^{(\ell)}\setminus


t^{(\ell)})\cup t^{(\ell)}-h$;


 $d^{(\ell+1)}=d^{(\ell)}-2(-


1)^{s^{(\ell)}}\gamma^{(\ell)}f_h(t^{(\ell)}-h)$; вычислим шаг


 $\sigma(t^{(\ell)}-h)=-\delta_h(t^{(\ell)}-


2h)/\Delta\delta_h(t^{(\ell)}-2h)$. Вместо значения


 $G(t^{(\ell)})$ в память ЭВМ занесем $G(t^{(\ell)}-h)$.


\item[A.3.] Если $t^{(\ell)}=t_*$, то


$T_{\nop0}^{(\ell+1)}=T_{\nop0}^{(\ell)}\cup t^{(\ell)}+h$


 (в $T_{\nop0}^{(\ell)}$ занесем новый неопорный нуль $t^{(\ell)}+h$);


положим $d^{(\ell+1)}=d^{(\ell)}+2\gamma^{(\ell)}f_h(t^{(\ell)})$,


$\gamma^{(\ell+1)}=-\gamma^{(\ell)}$,


 $\vartheta(t^{(\ell)})=1$, ${s^*}^{(\ell+1)}={s^*}^{(\ell)}+1$,


перенумеруем


 точки множества $T_{\opn}$. Занесем в память ЭВМ $G(t^{(\ell)}+h)$ и


найдем новый


 шаг $\sigma(t^{(\ell)}+h)=-


\delta_h(t^{(\ell)}+h)/\Delta\delta_h(t^{(\ell)}+h)$.


Эта ситуация трактуется как появление нового нуля функции


(\ref{vozm-koupravlen}) на


левом конце $T_h$.


\item[A.4.] Если $t^{(\ell)}=t^*$, то


$T_{\nop0}^{(\ell+1)}=T_{\nop0}^{(\ell)}\cup t^{(\ell)}-h$


(в $T_{\nop0}^{(\ell)}$ занесем новый неопорный нуль $t^{(\ell)}-h$);


положим $d^{(\ell+1)}=d^{(\ell)}+2(-


1)^{{s^*}^{(\ell)}}\gamma^{(\ell)}f_h(t^{(\ell)}-h)$,


$\vartheta(t^{(\ell)})=-1$, ${s^*}^{(\ell+1)}={s^*}^{(\ell)}+1$,


перенумеруем


точки множества $T_{\opn}$. Занесем в память ЭВМ $G(t^{(\ell)}-h)$ и


найдем шаг $\sigma(t^{(\ell)}-h)=-\delta_h(t^{(\ell)}-


2h)/\Delta\delta_h(t^{(\ell)}-2h)$.


Ситуация трактуется как возникновение новoго нуля функции


(\ref{vozm-koupravlen}) на правом конце $T_h$.


\end{itemize}





\noindent{}B. Пусть шаг в формуле (\ref{el-step}) достигся в два момента


$t^{(\ell)}$ и $t^{(\ell)}+h$:


 $\sigma(t^{(\ell)})=\sigma(t^{(\ell)}+h)$.


\begin{itemize}


\item[B.1.] Если $t^{(\ell)}=t_*$, то


$T_{\nop0}^{(\ell+1)}=T_{\nop0}^{(\ell)}\setminus t^{(\ell)}+h$


 (из множества $T_{\nop0}^{(\ell)}$ удалим точку $t^{(\ell)}+h$);


положим


 $d^{(\ell+1)}=d^{(\ell)}+2\gamma^{(\ell)}f_h(t^{(\ell)})$,


$\gamma^{(\ell+1)}=-\gamma^{(\ell)}$,


 $\sigma(t^{(\ell)})=\infty$, ${s^*}^{(\ell+1)}={s^*}^{(\ell)}-1$,


 перенумеруем точки множества $T_{\opn}$. Из памяти ЭВМ удалим значение


$G(t^{(\ell)}+h)$.


 В этом случае нуль $t^{(\ell)}+h$ переместится в точку $t_*$, если


$t_*\not\in T_{\op}$, или в точку


 $t_*+h$, если $t_*\in T_{\op}$. При $\sigma>\sigma^{(\ell)}$ эта


 ситуация трактуется как исчезновение нуля $t^{(\ell)}+h$ через левую


границу $T_h$.


\item[B.2.] Если $t^{(\ell)}+h=t^*$, то


$T_{\nop0}^{(\ell+1)}=T_{\nop0}^{(\ell)}\setminus t^{(\ell)}$


 (из $T_{\nop0}^{(\ell)}$ удалим точку $t^{(\ell)}$); положим


 $d^{(\ell+1)}=d^{(\ell)}+2(-


1)^{{s^*}^{(\ell)}}\gamma^{(\ell)}f_h(t^{(\ell)})$;


 $\sigma(t^{(\ell)}+h)=\infty$; ${s^*}^{(\ell+1)}={s^*}^{(\ell)}-1$,


 перенумеруем точки множества $T_{\opn}$. Из памяти ЭВМ удалим


$G(t^{(\ell)})$.


 При $\sigma>\sigma^{(\ell)}$ ситуация трактуется как исчезновение нуля


$t^{(\ell)}$ через


 правую границу множества $T_h$.


\item[B.3.] Если $t^{(\ell)},t^{(\ell)}+h{\in}


T_{\nop0}^{(\ell)}$, то положим


$T_{\nop0}^{(\ell+1)}{=}T_{\nop0}^{(\ell)}\setminus\{t^{(\ell)},\,t^{(\ell)


}+h\}$, $d^{(\ell+1)}=d^{(\ell)}+2(-


1)^{s^{(\ell)}}\gamma^{(\ell)}f_h(t^{(\ell)})$,


 ${s^*}^{(\ell+1)}={s^*}^{(\ell)}-2$, перенумеруем точки множества


$T_{\opn}$. Из


 памяти удалим $G(t^{(\ell)})$, $G(t^{(\ell)}+h)$. Эта ситуация


трактуется как слипание и


 исчезновение двух нулей в точке $t^{(\ell)}$ при


$\sigma>\sigma^{(\ell)}$.


 \end{itemize}





\noindent{}C. Если реализовались условия случая A и либо


1)~$t^{(\ell)}+h\in T_{\op}$ и $\vartheta(t^{(\ell)})=1$,


либо 2)~$t^{(\ell)}-2h\in T_{\op}$ и $\vartheta(t^{(\ell)})=-1$, то


ситуацию трактуем как переход подвижного нуля через опорный.


 \begin{itemize}


\item[C.1.] Положим


$T_{\nop0}^{(\ell+1)}=(T_{\nop0}^{(\ell)}\setminus t^{(\ell)})\cup


 t^{(\ell)}+2h$, $d^{(\ell+1)}=d^{(\ell)}+2(-


1)^{s^{(\ell)}}\gamma^{(\ell)}f_h(t^{(\ell)})$,


 вычислим новый шаг $\sigma(t^{(\ell)}+2h)=-


\delta_h(t^{(\ell)}+2h)/\Delta\delta_h(t^{(\ell)}+2h)$.


 Вместо значения $G(t^{(\ell)})$ в память ЭВМ занесем


$G(t^{(\ell)}+2h)$. Перенумеруем точки


 множества $T_{\opn}$.


\item[C.2.] Положим


$T_{\nop0}^{(\ell+1)}=(T_{\nop0}^{(\ell)}\setminus t^{(\ell)})\cup


 t^{(\ell)}-2h$, $d^{(\ell+1)}=d^{(\ell)}-2(-


1)^{s^{(\ell)}}\gamma^{(\ell)}f_h(t^{(\ell)}-h)$,


 вычислим новый шаг $\sigma(t^{(\ell)}-2h)=-\delta_h(t^{(\ell)}-


3h)/\Delta\delta_h(t^{(\ell)}-3h)$.


 Вместо значения $G(t^{(\ell)})$ в память ЭВМ занесем $G(t^{(\ell)}-


2h)$. Перенумеруем точки множества $T_{\opn}$.


\end{itemize}


Остальная информация переписывается для $(\ell+1)$-го шага без изменений.





{\it Заключительный шаг} (а). Информация для новой итерации преобразится


следующим образом.


\begin{itemize}


 \item[A.1.] $\ov{T}_{\op}=(T_{\op}\setminus t_0)\cup


 t^{(\ell)}$, $\ov{T}_{\nop0}=T_{\nop0}^{(\ell)}\setminus t^{(\ell)}$,


 $\ov{d}=d^{(\ell)}+(-1)^{s^{(\ell)}}\gamma^{(\ell)}f_h(t^{(\ell)})$.


\item[A.2.] $\ov{T}_{\op}=(T_{\op}\setminus t_0)\cup


 t^{(\ell)}-h$, $\ov{T}_{\nop0}=T_{\nop0}^{(\ell)}\setminus


t^{(\ell)}$,


 $\ov{d}=d^{(\ell)}-(-1)^{s^{(\ell)}}\gamma^{(\ell)}f_h(t^{(\ell)}-h)$.


 Вместо значения $G(t^{(\ell)})$ в память ЭВМ занесем


 $G(t^{(\ell)}-h)$.


\item[A.3.] $\ov{T}_{\op}=(T_{\op}\setminus t_0)\cup


 t^{(\ell)}$, $\ov{T}_{\nop0}=T_{\nop0}^{(\ell)}$,


 $\ov{d}=d^{(\ell)}+\gamma^{(\ell)}f_h(t^{(\ell)})$.


\item[A.4.] $\ov{T}_{\op}=(T_{\op}\setminus t_0)\cup


 t^{(\ell)}-h$, $\ov{T}_{\nop0}=T_{\nop0}^{(\ell)}$,


 $\ov{d}=d^{(\ell)}+(-1)^{{s^*}^{(\ell)}}\gamma^{(\ell)}f_h(t^{(\ell)}-


h)$, $\ov{s}^*={s^*}^{(\ell)}+1$.


 Перенумеруем точки множества $T_{\opn}$. Запомним $G(t^{(\ell)}-h)$.


\item[B.1.] $\ov{T}_{\op}=(T_{\op}\setminus t_0)\cup


 t^{(\ell)}$, $\ov{T}_{\nop0}=T_{\nop0}^{(\ell)}\setminus


t^{(\ell)}+h$,


 $\ov{d}=d^{(\ell)}+\gamma^{(\ell)}f_h(t^{(\ell)})$, $\ov{\gamma}=-


\gamma^{(\ell)}$,


 $\ov{s}^*={s^*}^{(\ell)}-1$. Перенумеруем точки множества $T_{\opn}$.


 Из памяти удалим значение $G(t^{(\ell)}+h)$.


\item[B.2.] $\ov{T}_{\op}=(T_{\op}\setminus t_0)\cup


 t^{(\ell)}$, $\ov{T}_{\nop0}=T_{\nop0}^{(\ell)}\setminus t^{(\ell)}$,


 $\ov{d}=d^{(\ell)}+(-


1)^{{s^*}^{(\ell)}}\gamma^{(\ell)}f_h(t^{(\ell)})$.


\end{itemize}


Новая опора имеет вид $\ov{K}_{\op}=\{I_{\op}, \ov{T}_{\op}\}$. Заменив


в матрице $P_{|\op|}$ столбец $f_h(t_0)$


на $f_h(t^{(\ell)}-h)$ в случаях A.2 или A.4 либо на столбец


$f_h(t^{(\ell)})$ в остальных случаях, получим


$\ov{P}_{|\op|}$.





\begin{notenonum}


Если $t_*\in T_{\op}$, то в случаях A.3 и B.1 вместо


точки $t_*$ рассмотрим $t_*+h$ и


подсчитаем $\sigma(t_*+h)=-\delta_h(t_*+h)/\Delta\delta_h(t_*+h)$.


Аналогично, если $t^*-h\in T_{\op}$, то вместо


точки $t^*$ рассмотрим $t^*-h$ и найдем $\sigma(t^*-h)=-\delta_h(t^*-


2h)/\Delta\delta_h(t^*-2h)$. При формировании


множества $T^{(\ell)}$ (см. подготовку информации 1)~--~7) перед $\ell$-м


шагом) произведем соответствующие замены.


\end{notenonum}





Выше ради упрощения изложения был опущен случай, когда в процессе


движения вдоль направления $\Delta\nu$


возникают новые нули внутри множества $T_h$. Для исследования этого случая


нужно дополнить информацию


(\ref{information}) совокупностью моментов $t$ и чисел $\sigma(t)$:


а)~точки $t\in T_h$ --- стационарные точки


варьированного коуправления (\ref{var-koupravlen}), т.\,е. при некотором


$\sigma>0$ выполняется


1)~$\delta_h(t,\sigma)>0$, $\delta_h(t-h,\sigma)>\delta_h(t,\sigma)$,


$\delta_h(t+h,\sigma)>\delta_h(t,\sigma)$ или


2)~$\delta_h(t,\sigma)<0$, $\delta_h(t-h,\sigma)<\delta_h(t,\sigma)$,


$\delta_h(t+h,\sigma)<\delta_h(t,\sigma)$;


б)~числа $\sigma(t)$ характеризуют величину шагов вдоль $\Delta\nu$, при


которых варьированное коуправление в момент


$t$ обращается в нуль: $\delta_h(t,\sigma(t))=0$. Следя за движением и


положением стационарных точек варьированного


коуправления, зафиксируем появление нового ``внутреннего'' нуля


коуправления. В сложных ситуациях могут


появиться и новые стационарные точки. Они обнаруживаются с помощью


стационарных точек первой производной


коуправления. Объем дополнительной информации зависит от сложности


конкретной задачи. Преобразование дополнительной


информации осуществляется по аналогии с описанным выше для стандартного


случая.





б)~Когда $\sigma^{(\ell)}=\sigma_{i^{(\ell)}}$, то


$\sigma_{i^{(\ell)}}=\infty$. Затем перейдем к $(\ell+1)$-му шагу.





{\it Заключительный шаг} (б). Сформируем информацию для следующего шага:


$\ov{T}_{\op}=T_{\op}\setminus t_0$,


$\ov{I}_{\op}=I_{\op}\setminus i^{(\ell)}$,


$\ov{K}_{\op}=\{\ov{I}_{\op},\,\ov{T}_{\op}\}$,


$\ov{T}_{\nop0}=T_{\nop}^{(\ell)}$, $\ov{d}=d^{(\ell)}$,


$\ov{\gamma}=\gamma^{(\ell)}$,


$\ov{s}^*={s^*}^{(\ell)}$, $G(t)$, $t\in \ov{T}_{\nop0}$. Матрицу


$\ov{P}_{|\op|}$ получим, удалив из


$P_{|\op|}$ столбец $f_h(t_0)$.





Согласно [16] обязательно найдется такое $\ell_0$, что


$\eta^{(\ell_0)}>0$, $\eta^{(\ell_0+1)}\leq0$.





Ситуация 1) ($\rho_0=\rho(t_0)$) на с.\,\pageref{rho} исследована до конца.





2)~($\rho_0=\rho_{i_0}$). Направление изменения вектора потенциалов


$\Delta\nu$ найдем из уравнения


\begin{gather*}


-P'_{\op}\Delta\nu=({f_h}_{i_0}(t),\,t\in


T_{\op};\,{f_l}_{i_0})\mu\\


(\mu=1,\text{ если }\zeta_{i_0}(t^*)<{g_*}_{i_0}(\alpha);\ \ 


\mu=-1,\text{ если }\zeta_{i_0}(t^*)>g^*_{i_0}(\alpha)).


\end{gather*}





Направление изменения коуправления построим по формуле


(\ref{var-koupravlen}). Начальная скорость изменения


двойственного критерия качества вдоль $\Delta\nu$ теперь равна


$\eta^{(1)}=\rho_{i_0}>0$.





Операции по вычислению ``длинного'' шага аналогичны описанным для


ситуации~1), но ``нулевой'' шаг здесь не


проводится. Новая опора строится по следующим правилам.





а)~$(\sigma^*=\sigma(t^{(\ell_0)}))$


$\ov{K}_{\op}=\{\ov{I}_{\op},\ov{T}_{\op}\}$,


$\ov{I}_{\op}=I_{\op}\cup


i_0$. В случаях A.2 или A.4 $\ov{T}_{\op}=T_{\op}\cup t^{(\ell_0)}-h$,


в матрицу $\ov{P}_{|\op|}$ добавим новый


столбец $f_h(t^{(\ell_0)}-h)$. В остальных случаях


$\ov{T}_{\op}=T_{\op}\cup t^{(\ell_0)}$, в матрицу


$\ov{P}_{|\op|}$ добавим новый столбец $f_h(t^{(\ell_0)})$. Правила


построения множества $\ov{T}_{\nop0}$


аналогичны случаю 1а).





б)~$(\sigma^*=\sigma_{i^{(\ell_0)}})$


$\ov{K}_{\op}=\{\ov{I}_{\op},T_{\op}\}$,


$\ov{I}_{\op}=(I_{\op}\setminus i^{(\ell_0)})\cup i_0$,


$\ov{T}_{\nop0}=T_{\nop0}^{(\ell_0)}$,


$\ov{P}_{|\op|}=P_{|\op|}$.





Двойственный метод завершит работу построением оптимальной опоры и


сопровождающего ее псевдоуправления, которое


будет оптимальной программой линейной задачи


(\ref{vspom-problem}).








\section{Доводка}





Пусть $\wt{u}(t)$, $t\in T$, --- оптимальная программа


вспомогательной задачи (\ref{vspom-problem}), удовлетворяющая


условию перехода к процедуре доводки; $\wt{t}_1,


\wt{t}_2,\dotsc, \wt{t}_{p-1}$


--- точки переключения программы.


Согласно [17] для оптимальности программы задачи (\ref{problem})


необходимо и достаточно, чтобы она имела вид


\begin{equation}\label{krit-optim-nepr}


u(t)=-\sign\delta(t),\ \ t\in T,


\end{equation}


где $\delta(t)=\psi'(t)b(t)$, $t\in T$, --- коуправление,


$\psi(t)$, $t\in T$, --- котраектория


--- решение сопряженного уравнения (\ref{sopr-sistem}) с начальным


условием $\psi(t^*)=\partial\varphi(x(t^*))/\partial x$.





Отсюда следует, что моменты переключения $t_1^0, t_2^0,\dotsc,


t_{p-1}^0$ оптимальной программы являются решением уравнений


\begin{equation}\label{uravn-dododki}


\delta(t_j)=0,\ \ j=\overline{1,p-1}.


\end{equation}





В подробной записи уравнения (\ref{uravn-dododki}) имеют вид


$\psi'(t^*)F(t^*,t_j)b(t_j)=0$, $j=\overline{1,p-1}$.





Назовем систему уравнений (\ref{uravn-dododki}) уравнениями


доводки. При $\dot{\delta}(t_j)\neq0$, $j=\overline{1,p-1}$,


матрица Якоби системы (\ref{uravn-dododki}) неособая. В этом


случае уравнения доводки можно решить методом Ньютона. В качестве


начального приближения возьмем моменты $\wt{t}_j$,


$j=\overline{1,p-1}$, полученные после решения вспомогательной


задачи (\ref{vspom-problem}). Если число $\varepsilon>0$ выбрано


достаточно малым, то метод Ньютона будет сходиться с квадратичной


скоростью. Быстрая реализация метода Ньютона для систем вида


(\ref{uravn-dododki}) описана в [17]. В результате решения


уравнений (\ref{uravn-dododki}) будет построена оптимальная


программа задачи (\ref{problem}) в классе кусочно-непрерывных


функций.





Для приложений, в которых управляющие воздействия реализуются с


помощью устройств дискретного действия, большую роль играют


оптимальные программы в классе дискретных функций. Такие программы


можно построить с помощью дискретной процедуры доводки. В классе


дискретных программ критерий оптимальности (\ref{krit-optim-nepr})


примет вид


\begin{equation}\label{krit-optim-diskret}


u_h(t)=-\sign\delta_h(t),\ \ t\in T_h,


\end{equation}


где $\delta_h(t)=\int\limits_t^{t+h}\delta(\tau)d\tau$,


$\delta(t)=\psi'_h(t)b(t)$; $\psi_h(t)$, $t\in T$, --- решение


уравнения (\ref{sopr-sistem}) с начальным условием


$\psi_h(t^*)=\partial\varphi(x_h(t^*))/\partial x$; $x_h(t)$,


$t\in T$, --- траектория, порожденная дискретной программой


(\ref{krit-optim-diskret}).





Используя этот критерий, можно записать дискретный аналог


уравнений доводки (\ref{uravn-dododki}). В отличие от непрерывного


случая неизвестными в этих уравнениях будут значения оптимальной


программы в моменты переключения.





В случае малых периодов квантования $h$ при построении оптимальных


дискретных программ можно ограничиться процедурой преддоводки,


которая состоит в следующем. По оптимальной кусочно-непрерывной


программе $u^0(t)$, $t\in T$, найдем интервалы $[t_i,t_i+h]$,


$i=\overline{1,p-1}$, содержащие точки переключения программы


$t_i^0$, $i=\overline{1,p-1}$. Положим $u_h^0(t)=u^0(t)$, $t\in


T\setminus\bigcup\limits_{i=1}^{p-1}[t_i,t_i+h[$. Значения $u^0(t)$,


$t\in [t_i,t_i+h[$, $i=\overline{1,p-1}$, заменим на


$u_h^0(t_i)=2u^0(t_i)(t_i^0-t_i-h/2)$, $i=\overline{1, p-1}$.








\section{Синтез оптимальных управлений типа обратной связи}





Как известно, оптимальные программы позволяют лишь выявлять


потенциальные возможности систем управления. В реальных процессах


используются, как правило, управления типа обратной


связи. Определим ОУ типа (дискретной) обратной связи (с периодом


реализации $h_0=kh$, $k\geq1$). С этой целью погрузим задачу


(\ref{problem}) в семейство задач


\begin{equation}\label{fam-problems}


\varphi(x(t^*))\rightarrow\min,\ \ \Dot{x}=A(t)x+b(t)u,\ \ 


x(\tau)=z,\ \ |u(t)|\leq1,\ \ t\in T(\tau)=[\tau,t^*],


\end{equation}


зависящее от скаляра $\tau\in T_{h_0}$ и $n$-вектора $z$.





Пусть $u^0(t\mid \tau,z)$, $t\in T(\tau)$, --- оптимальная программа


задачи (\ref{fam-problems}) для позиции $(\tau,z)$. Следуя теории


оптимальных процессов, функцию


\begin{equation}\label{position-sol}


u^0(\tau,z)=u^0(\tau\mid\tau,z),\ \ \tau\in T_{h_0},\ \ z\in R^n,


\end{equation}


назовем ОУ типа (дискретной) обратной связи (позиционным решением


задачи (\ref{problem})), построение функции (\ref{position-sol})


--- синтезом оптимальной обратной связи (синтезом оптимальной


системы).





Использование обратных связей (\ref{position-sol}) предполагает,


что в процессе управления состояния системы (\ref{problem})


измеряются в дискретные моменты $t\in T_{h_0}$ и по этой


информации вырабатываются управляющие воздействия.





Замену в (\ref{problem}) управления $u$ на функцию


(\ref{position-sol}) называют замыканием системы управления.


Траектория замкнутой системы


\begin{equation}\label{close-sistem}


\Dot{x}=A(t)x+b(t)u^0(t,x)+w(t),\,x(t_*)=x_0,


\end{equation}


при постоянно действующем кусочно-непрерывном возмущении $w(t)$,


$t\in T$, представляет решение уравнения


\begin{gather*}


\Dot{x}=A(t)x+b(t)u(t)+w(t),\ \ x(t_*)=x_0,\\


u(t)=u^0(\tau,x(\tau)),\ \ t\in


[\tau,\tau+h_0[,\ \ \tau=t_*+sh_0,\ \ s=\overline{0,N-1}.


\end{gather*}





Будем считать, что уравнение (\ref{close-sistem}) описывает


поведение физического прототипа математической модели


(\ref{problem}). В этом случае функция $w(t)$, $t\in T$, содержит


неточности математического моделирования и постоянно действующие


возмущения.





Пусть по ходу некоторого конкретного процесса управления


реализуется возмущение $w^*(t)$, $t\in T$. Ему будет


соответствовать траектория $x^*(t)$, $t\in T$, замкнутой системы


(\ref{close-sistem}), удовлетворяющая тождеству


\begin{equation}\label{identity}


\Dot{x}^*(t)\equiv A(t)x^*(t)+b(t)u^0(t,x^*(t))+w^*(t),\ \ t\in T.


\end{equation}





Из тождества (\ref{identity}) видно, что в процессе управления


обратная связь (\ref{position-sol}) используется не полностью (не


для всех $z\in R^n$). Нужны лишь ее значения


$u^*(t)=u^0(t,x^*(t))$, $t\in T_{h_0}$, вдоль изолированной


траектории $x^*(t)$, $t\in T$. При этом нет необходимости знать


функцию $u^*(t)$, $t\in T_{h_0}$, заранее, достаточно уметь


вычислять текущие значения $u^*(\tau)$ по мере поступления


измерений текущих состояний $x^*(\tau)$. Функцию $u^*(t)$, $t\in


T$, назовем реализацией оптимальной обратной связи в конкретном


процессе управления.





Будем говорить, что построение $u^*(t)$, $t\in T$, осуществляется


в режиме реального времени, если в каждой текущей позиции


$(\tau,x^*(\tau))$ время $s(\tau)$ вычисления значения


$u^*(\tau)=u^0(\tau,x^*(\tau))$ не превосходит $h_0>0$.


Устройство, способное выполнять такую работу, называется


оптимальным регулятором.





Алгоритм работы оптимального регулятора базируется на описанном


выше двойственном методе.





До начала процесса управления оптимальный регулятор вычисляет


оптимальную программу $u^0(t\mid t_*,x_0)$, $t\in T$, для начальной


позиции $(t_*,x_0)$. Время на выполнение этой работы не имеет


никакого значения. Для дальнейшего процесса управления


запоминается следующая информация: 1)~моменты переключения


$t_i^0=t_i^0(t_*,x_0)$, $i=\overline{1, p(t_*)-1}$; 2)~оптимальное


значение $\alpha^0=\alpha^0(t_*,x_0)$ критерия качества задачи


(\ref{problem}); 3)~оптимальная опора $K_{\op}^0(t_*,x_0)$


линейной задачи (\ref{vspom-problem}) для


$\alpha=\alpha^0(t_*,x_0)$; 4)~совокупность векторов аппроксимации


$r_i$, $i\in I(t_*)$, использованная при построении оптимальной


программы последней линейной задачи; 5)~начальный участок


оптимальной программы $u^0(t\mid t_*,x_0)$, $t\in [t_*,t_*+2h_0]$.





Процесс управления стартует в момент $t_*$, когда регулятор


начинает подавать на вход объекта управления управляющее


воздействие $u^*(t)=u^0(t|t_*,x_0)$, $t\geq t_*$.





Предположим, что процесс управления проведен на промежутке времени


$[t_*,\tau[$, в результате которого выработано управляющее


воздействие $u^*(t)$, $t\in [t_*,\tau[$. Пусть под действием этого


управления и реализовавшегося возмущения $w^*(t)$, $t\in


[t_*,\tau]$, объект управления в момент $\tau$ оказался в


состоянии $x^*(\tau)$, информация о котором поступила в регулятор.


Кроме этой информации оптимальный регулятор хранит в памяти


данные, полученные в предыдущий момент $\tau-h_0$: 1)~точки


переключения $t_i^0(\tau-h_0,x^*(\tau-h_0))$, $i=\overline{1,


p(\tau-h_0)-1}$; 2)~оптимальное значение


$\alpha^0(\tau-h_0,x^*(\tau-h_0))$ критерия качества задачи


(\ref{problem}); 3)~оптимальную опору


$K_{\op}^0(\tau-h_0,x^*(\tau-h_0))$ линейной


задачи (\ref{vspom-problem}) при


$\alpha=\alpha^0(\tau-h_0,x^*(\tau-h_0))$; 4)~совокупность


векторов аппроксимации $r_i$, $i\in I(\tau-h_0)$, использованную


при решении последней линейной задачи. Имея эту информацию,


оптимальный регулятор начинает параллельно осуществлять следующие


операции.





1.~Процедуру доводки с начальными значениями точек переключения


$t_i^0(\tau-h_0,x^*(\tau-h_0))$, $i=\overline{1, p(\tau-h_0)-1}$,


и начальным состоянием $x^*(\tau)$.





2.~Решение линейной задачи (\ref{vspom-problem}) с


$\alpha=\alpha^0(\tau-h_0,x^*(\tau-h_0))-\lambda$, начальной


опорой $K_{\op}^0(\tau-h_0,x^*(\tau-h_0))$ и векторами


аппроксимации $r_i$, $i\in I(\tau-h_0)$ ($\lambda>0$


--- параметр метода).





3.~Решение линейной задачи (\ref{vspom-problem}) с


$\alpha=\alpha^0(\tau-h_0,x^*(\tau-h_0))+\lambda$, начальной


опорой $K_{\op}^0(\tau-h_0,x^*(\tau-h_0))$ и векторами


аппроксимации $r_i$, $i\in I(\tau-h_0)$.





Если процедура доводки сходится, то по ее результатам информация


1)--4), использованная для момента $\tau$, обновляется для


момента $\tau+h_0$. В противном случае, решая линейные задачи


2, 3, добиваемся такой точности аппроксимации, при которой


станут выполняться условия перехода к процедуре доводки и


сходимость последней.





Управление $u^0(t\mid \tau,x^*(\tau))$, $t\in [\tau, \tau+2h_0]$,


подается на вход объекта управления, начиная с момента


$\tau+s^*(\tau)$, где $s^*(\tau)$ --- время, потраченное на


осуществление операций 1, 2, 3. Как показано в [15], [17],


описанные операции по реализации оптимальной обратной связи могут


быть с помощью существующих микропроцессоров осуществлены очень


быстро. Поэтому предложенные методы можно применять для


оптимального управления динамическими системами достаточно


высокого порядка. При этом задержки $s^*(\tau)$, $\tau\in


T_{h_0}$, которые вызваны выполнением указанных выше операций, не


оказывают существенного влияния на качество переходных процессов


[18].











\begin{thebibliography}{99}





\bibitem{Pontr}


Понтрягин Л.С., Болтянский В.Г., Гамкрелидзе Р.В., Мищенко Е.Ф.


{\it Математическая теория оптимальных


процессов}. -- М.: Наука, 1969. -- 384 c.





\bibitem{OLS}


Габасов Р., Кириллова Ф.М. {\it Оптимизация линейных систем. Методы


функционального анализа}. --


Минск: Изд-во Белорусск. ун-та, 1973. -- 246 c.





\bibitem{Mor}


Morari M., Jay H. Lee. {\it Model predictive control$:$ past, present and


future} // Comput. \& Chemical


Engineer. -- 1999. -- \No\,23(4/5). -- P.\,667--682.





\bibitem{Qin}


Qin S.J., Badgwell T.A. {\it An overview of industrial model predective


control technology} // Proc. Fifth


International Conference on Chemical Process Control --- CPC\,V / Eds.


J.K.\,Kantor, C.E.\,Garcia, and B.\,Carnohan.


American Institute of Chemical Engineers, 1996. -- P.\,232--256.





\bibitem{DAN1991}


Габасов Р., Кириллова Ф.М., Костюкова О.И. {\it Построение оптимальных


управлений типа обратной связи в линейной


задаче} // ДАН СССР. -- 1991. -- Т.\,320. -- \No\,6. -- С.\,1294--1299.





\bibitem{TK1992}


Габасов Р., Кириллова Ф.М., Костюкова О.И. {\it Оптимизация линейной


системы управления в режиме реального времени}


// Изв. РАН. Сер. техн. кибернет. -- 1992. -- \No\,4. -- С.\,3--19.





\bibitem{SIAM}


Gabasov R., Kirillova F.M., Balashevich N.V. {\it On the synthesis


problem for optimal control systems} // SIAM


Journal on Control and Optimization. -- 2001. -- V.\,39. --


\No\,4. -- P.\,1008--1042.





\bibitem{LubAiT1997}


Габасов Р., Лубочкин А.В. {\it Синтез регулятора для одной 


линейно-квадратичной задачи оптимального управления} //


Автоматика и телемеханика. -- 1997. -- \No\,9. -- С.\,3--14.





\bibitem{DmAiT2002}


Габасов Р., Дмитрук Н.М., Кириллова Ф.М. {\it Оптимизация многомерных


систем управления с параллелепипедными


ограничениями} // Автоматика и телемеханика. -- 2002. --


\No\,3. -- С.\,3--26.





\bibitem{TK1994}


Габасов Р., Кириллова Ф.М., Костюкова О.И. {\it К методам стабилизации


динамических систем} // Изв. РАН. Сер. техн.


кибернет. -- 1994. -- \No\,3. -- С.\,67--77.





\bibitem{RuzhDNANB2000}


Габасов Р., Кириллова Ф.М., Ружицкая Е.А. {\it Стабилизация перевернутого


маятника} // Докл. АН Беларуси. -- 2000. --


Т.\,44. -- \No\,2. -- С.\,9--12.





\bibitem{RuzhTiSU2001}


Габасов Р., Кириллова Ф.М., Ружицкая Е.А. {\it Демпфирование и


стабилизация маятника при больших начальных


возмущениях} // Изв. РАН. Теория и системы управления. -- 2001. --


\No\,1. -- С.\,29--38.





\bibitem{LubPMM1998}


Габасов Р., Лубочкин А.В. {\it Стабилизация линейных динамических систем


оптимальными управлениями


линейно-квадратичных задач} // ПММ. -- 1998. -- Т.\,62. --


Вып.\,4. -- С.\,556--565.





\bibitem{KMO2}


Габасов Р., Кириллова Ф.М. {\it Конструктивные методы оптимизации}. Ч.\,2.


{\em Задачи управления}. -- Минск: Изд-во Университетское,


1984. -- 207 c.





\bibitem{BalZHVM2000}


Балашевич Н.В., Габасов Р., Кириллова Ф.М. {\it Численные методы


программной и позиционной оптимизации линейных


систем управления} // Журн. вычисл. матем. и матем. физ. -- 2000. --


Т.\,40. -- \No\,6. -- С.\,838--859.





\bibitem{KMO1}


Габасов Р., Кириллова Ф.М., Тятюшкин А.И. {\it Конструктивные методы


оптимизации}. Ч.\,1. {\em Линейные задачи}. -- Минск: Изд-во Университетское,


1984. -- 124 c.





\bibitem{KMO4}


Габасов Р., Кириллова Ф.М., Костюкова О.И., Ракецкий В.М. {\it


Конструктивные методы оптимизации}. Ч.\,4. {\em Выпуклые


задачи}. -- Минск: Изд-во Университетское, 1987. -- 223 c.





\bibitem{BalZHVM2001}


Балашевич Н.В., Габасов Р., Кириллова Ф.М. {\it Алгоритмы программной и


позиционной оптимизации линейных систем


управления с промежуточными фазовыми ограничениями} // Журн. вычисл. матем. и


матем. физ. -- 2001. -- Т.\,41. -- \No\,10. -- С.\,1485--1504.





\end{thebibliography}





\vskip 2cm





\post{\it Белорусский государственный университет}{\it Поступила}


\smallskip


\post{\it Институт математики Национальной}{17.12.2003}


\post{\it Академии наук Беларуси}{}


\smallskip


\post{\it Брестский государственный}{}


\post{\it технический университет}{}





\label{end}


\end{document}














В данной работе при решении примеров оказались


достаточными следующие операции.





{\it Шаг~1.} Найдем решение $\beta^{(1)}$ задачи


(\ref{vspom-problem}) с


$\alpha=\alpha^{(1)}=\varphi(x_0+(1-\beta^*)l)$.





{\it Шаг~2.} В качестве $\alpha^{(2)}$ выберем на плоскости


$\alpha$O$\beta$ точку пересечения с осью O$\alpha$ прямой,


проведенной через точки $(\varphi^*; \beta^*)$ и $(\alpha^{(1)};


\beta^{(1)})$. Найдем $\beta^{(2)}$ и


$\frac{\partial\beta}{\partial\alpha}|_{\alpha=\alpha^{(2)}}$,


решив задачу (\ref{vspom-problem}) с $\alpha=\alpha^{(2)}$.





{\it Шаг~3.} Пусть $\alpha^{(3)}$ --- точка пересечения с осью


O$\alpha$ квадратичной параболы, проведенной через точки


$(\alpha^{(1)}; \beta^{(1)})$, $(\alpha^{(2)}; \beta^{(2)})$,


производная которой в точке $(\alpha^{(2)}; \beta^{(2)})$ равна


$\frac{\partial\beta}{\partial\alpha}|_{\alpha=\alpha^{(2)}}$.


Решив задачу (\ref{vspom-problem}) с $\alpha=\alpha^{(3)}$,


получим $\beta^{(3)}$,


$\frac{\partial\beta}{\partial\alpha}|_{\alpha=\alpha^{(3)}}$.





{\it Шаг~4.} Найдем точку пересечения $\alpha^{(4)}$ с осью


O$\alpha$ кубической параболы, проведенной через точки


$(\beta^{(2)}; \alpha^{(2)})$, $(\beta^{(3)}; \alpha^{(3)})$ и


имеющей в этих точках производные


$\frac{\partial\beta}{\partial\alpha}|_{\alpha=\alpha^{(2)}}$,


$\frac{\partial\beta}{\partial\alpha}|_{\alpha=\alpha^{(3)}}$.


Решив задачу (\ref{vspom-problem}) с $\alpha=\alpha^{(4)}$, найдем


$\beta^{(4)}$.





Полученное значение $\beta^{(4)}$ во всех решенных примерах было


равно нулю с высокой точностью. Поэтому полагали


$\wt{\alpha}^0=\alpha^{(4)}$.





Вспомогательную линейную задачу (\ref{vspom-problem}) с


обновленным множеством векторов аппроксимации начнем решать с


параметром $\alpha^{(1)}$, равным оптимальному значению


$\wt{\alpha}^0$ критерия качества задачи


(\ref{approxim-problem}) с предыдущим набором векторов


аппроксимации. Решив задачу (\ref{vspom-problem}) с


$\alpha=\alpha^{(1)}$, получим $\beta^{(1)}$,


$\frac{\partial\beta}{\partial\alpha}|_{\alpha=\alpha^{(1)}}$.


Построив прямую по $\alpha^{(1)}$, $\beta^{(1)}$,


$\frac{\partial\beta}{\partial\alpha}|_{\alpha=\alpha^{(1)}}$,


найдем точку пересечения $\alpha^{(2)}$ прямой с осью O$\alpha$.


Затем вычисления продолжим с третьего шага по описанной выше


схеме.


